\section{Introduction}
\label{sec:intro}

Surgeons learn by being assessed, and suturing, passing a curved needle
through tissue with an instrument, is among the skills assessed most. An
expert watching a trainee can tell a good stitch from a poor one, but such
review takes time, varies between reviewers, and ``poor'' alone does not say
what to change. EASE~\cite{Haque2022} is a rating tool for robotic suturing
that assesses each sub-skill against a defined criterion, such as where along
the needle it is held or at what angle, so that a score points to what to
correct. Rating every stitch by hand, however, takes expert time that
training rarely has to spare.

Automated assessment from surgical video has moved from classifying
expertise on JIGSAWS~\cite{Funke2019,Liu2021} to clinically defined rating
tools~\cite{Kiyasseh2023}; for EASE, models trained on annotated stitches
predict each sub-skill as good or poor~\cite{Cui2024,Deo2026}, and a
general-purpose vision--language model (VLM) can be asked for a score
directly. In both cases the score comes out of a black box: attention maps
and evidence frames show where the model
looked~\cite{IsmailFawaz2019,Kiyasseh2023b}, but not how the score was
reached. Instrument keypoints and
masks can be detected in surgical video~\cite{Ghanekar2025}, and the 3D pose
of a suture needle has been estimated, but with stereo cameras, calibration,
or trained detectors~\cite{Chiu2022,Li2024}.

Our observation is that EASE defines its criteria geometrically: the grasp
position along the needle, the angle between the needle and the jaws, and
the depth of the needle in the jaws. We therefore measure rather than
predict: we measure what the criteria define and score the measurements by
the criteria (Fig.~\ref{fig:overview}). Events detected
in the video divide it into stitches and into the window that each sub-skill
refers to; within it, per-frame keypoints and a needle mask give the
geometry, which is measured in 3D from the monocular video and scored by the
EASE criteria. Every score thus comes with the frames and the value behind
it, and can be verified.

Our contributions are threefold. (1) A measurement-based formulation of EASE
scoring: a score is a geometric quantity compared with the EASE criterion,
measured within the event-defined window in which the hold does not change,
so that every frame in it is a valid observation; the perception that
provides the keypoints and masks is replaceable, and this paper implements
it without training any model, with a general-purpose VLM and a few
examples.
(2) Geometric measurement from monocular video without calibration: we fit
one 3D arc to the needle across the whole video to locate the grasp on it,
lift the needle and the jaws to 3D for the hold angle using the fact that
projection only shortens, and read the hold depth as a ratio along the jaws,
which projection preserves. (3) Validation on JIGSAWS~\cite{Gao2014} with surgeons split
between development and test: run fully automatically from the raw video,
the method reproduces a surgeon's EASE scores at least as well as direct
grading by a state-of-the-art VLM given favorable conditions, and most of
its errors lie near a threshold.
